\titledquestion{Implement factor graph SLAM}
\label{q:task01}
A factor graph SLAM system has three jobs: build the graph, recover a covariance from it, and use that covariance to decide which measurement belongs to which landmark.
The back-end (iSAM2, through GTSAM) does the optimization for you; everything below is the part it does \emph{not} do.

All functions are marked with \pythoninline{# TODO} in the code. Run the tests with \pythoninline{pytest}: how many of them pass is part of how we evaluate your submission as a whole. The GTSAM calls you need are listed in \pythoninline{docs/GTSAM.md}, and your editor shows their signatures.

Two conventions are worth fixing in your head before you start, because both produce systems that run happily while being wrong:
\begin{itemize}
    \item Measurements are ordered \([\text{range}, \text{bearing}]\) everywhere in this code base. GTSAM's \pythoninline{BearingRangeFactor2D} takes \emph{bearing first}.
    \item The covariance GTSAM reports for a \pythoninline{Pose2}, and the Jacobians returned by \pythoninline{Pose2.range} and \pythoninline{Pose2.bearing}, both live in the \emph{tangent space} at the current estimate (Sec.\ 6.2 in the book), not in raw \((x, y, \psi)\) coordinates. Used together they are consistent. Mixing either with a hand-derived \(\partial/\partial[x, y, \psi]\) is not.
\end{itemize}

\begin{parts}

    % ------------------------------------------------------------------ (a) % TODO ADD: gtsam.Pose2.Expmap expects np.array 
    \part Implement \pythoninline{preprocessing.relative_pose}\unskip
, which turns one wheel-encoder reading into a relative pose increment.

    The Victoria Park vehicle reports a forward speed \(v_e\) measured at the left rear wheel and a steering angle \(\alpha\). The encoder sits a distance \(H\) off the centre line, so during a turn it travels along a different arc than the centre of the rear axle. With \(L\) the axle distance, the kinematic bicycle model gives the body twist
    \begin{align*}
        v_b    & = \frac{v_e}{1 - \frac{H}{L}\tan\alpha}, &
        \omega & = \frac{v_b}{L}\tan\alpha,               &
        \xi    & = \begin{bmatrix} v_b & 0 & \omega\end{bmatrix}^T.
    \end{align*}
    Held constant over the interval, this twist integrates \emph{exactly} to \(\Delta T = \mathrm{Exp}(\xi\,\Delta t)\), which follows a circular arc.

    \begin{longhint}
        Do not Euler-integrate into \(\begin{bmatrix} v_b\Delta t & 0 & \omega\Delta t\end{bmatrix}\) and build a \pythoninline{Pose2} from that. The difference is second order in \(\omega \Delta t\): invisible in one step, and clearly visible after a few thousand of them. The exponential map is covered in Sec.\ 6.2.3; \pythoninline{gtsam.Pose2.Expmap} implements it, taking \(\xi\Delta t\) as a numpy array. The vehicle geometry is the \pythoninline{car} argument: \(H\) is \pythoninline{car.H} and \(L\) is \pythoninline{car.L}.
    \end{longhint}
    \begin{solution} \begin{pythoncode}
def relative_pose(vel_encoder: float, steer: float, dt: float, car: Car = CAR) -> gtsam.Pose2:
    tangent_steer = np.tan(steer)
    velocity_body = vel_encoder / (1.0 - (car.H / car.L) * tangent_steer)
    yaw_rate = (velocity_body / car.L) * tangent_steer

    twist_body = np.array([velocity_body, 0.0, yaw_rate], dtype=float)

    return gtsam.Pose2.Expmap(twist_body * dt)
\end{pythoncode}
 \end{solution}

    % ------------------------------------------------------------------ (b)
    \breakans
    \part Implement \pythoninline{preprocessing.preintegrate}\unskip
, which compounds a sequence of relative poses into a single one, with covariance.

    Wheel odometry arrives far faster than lidar scans. Rather than putting a pose in the graph for every encoder tick, all increments between two scans are compounded into one relative pose and one covariance, which enter the graph as a single \pythoninline{BetweenFactorPose2}. This is \emph{preintegration}, mentioned in Sec.\ 9.1 -- the same idea used for IMU data, in its simplest form.

    The mean is repeated composition. The covariance is the interesting part: composition is nonlinear, so each increment's covariance must be pushed through the Jacobians of the composition,
    \[
        \Sigma \leftarrow H_1 \Sigma H_1^T + H_2 \Sigma_i H_2^T,
    \]
    where \(H_1\) and \(H_2\) are the Jacobians of \pythoninline{compose} with respect to its first and second argument. This is ordinary linear covariance propagation; it just happens on a manifold, so the Jacobians are those of the group operation rather than of a vector sum.

    \begin{longhint}
        Start the recursion from ``no motion, known exactly'': the identity pose \pythoninline{gtsam.Pose2()} and a zero \(3\times3\) covariance. That is also the right answer for an empty list of increments. \pythoninline{gtsam.Pose2.compose} fills in both Jacobians for you if you pass two preallocated \pythoninline{np.zeros((3, 3), order="F")} arrays.
    \end{longhint}
    \begin{solution} \begin{pythoncode}
def preintegrate(
    poses: list[gtsam.Pose2],
    covariances: list[np.ndarray],
) -> tuple[gtsam.Pose2, np.ndarray]:
    compounded = gtsam.Pose2.Identity()
    compounded_cov = np.zeros((3, 3))

    for increment, increment_cov in zip(poses, covariances):
        H1 = np.zeros((3, 3), order="F")
        H2 = np.zeros((3, 3), order="F")

        compounded = compounded.compose(increment, H1, H2)
        compounded_cov = H1 @ compounded_cov @ H1.T + H2 @ increment_cov @ H2.T

    return compounded, compounded_cov
\end{pythoncode}
 \end{solution}

    % ------------------------------------------------------------------ (c)
    \breakans
    \part Build the graph. Each factor you add is literally one term of the negative log posterior (9.4), and therefore one term of the cost function (9.6) that the back-end minimizes.
    \begin{subparts}
        \subpart Implement \pythoninline{factor_graph.predict_pose}\unskip
, which returns the predicted pose
        \[
            \hat{x}_k = \hat{x}_{k-1} \oplus u_k :
        \]
        the previous pose estimate compounded with the odometry increment \(u_k\), which is expressed in the body frame of \(x_{k-1}\). In words: rotate the increment's translation by the previous heading and add it to the previous position, and add the two headings. The compounding \(\oplus\) is the inverse of the \(\ominus\) in (9.5); for a \pythoninline{gtsam.Pose2} it is \pythoninline{compose}.

        The prediction is only the starting value for the new variable \(x_k\): (9.13) is solved iteratively, so every variable needs one. It adds no information to the graph; the odometry factor in the next part does that.
        \begin{solution} \begin{pythoncode}
def predict_pose(previous_pose: gtsam.Pose2, relative_pose: gtsam.Pose2) -> gtsam.Pose2:
    return previous_pose.compose(relative_pose)
\end{pythoncode}
 \end{solution}

        \subpart Implement \pythoninline{factor_graph.add_odometry_factor}\unskip
, adding the middle term of (9.6) as a \pythoninline{BetweenFactorPose2}.
        \begin{longhint}
            The factor is constructed as
            \begin{center}\pythoninline{gtsam.BetweenFactorPose2(key_previous, key_current, measured, noise_model)}\end{center}
            and \pythoninline{graph.add(factor)} adds it to the graph. Use \pythoninline{gtsam.noiseModel.Gaussian.Covariance} and not \pythoninline{.Sigmas}: the covariance you preintegrated in (b) is not diagonal, and discarding its off-diagonal terms is an easy way to end up over-confident and inconsistent.
        \end{longhint}
        \begin{solution} \begin{pythoncode}
def add_odometry_factor(
    graph: gtsam.NonlinearFactorGraph,
    key_previous: int,
    key_current: int,
    relative_pose: gtsam.Pose2,
    relative_pose_cov: np.ndarray,
) -> None:
    noise = gtsam.noiseModel.Gaussian.Covariance(relative_pose_cov)
    graph.add(
        gtsam.BetweenFactorPose2(
            key_previous,
            key_current,
            relative_pose,
            noise,
        )
    )
\end{pythoncode}
 \end{solution}

        \subpart Implement \pythoninline{factor_graph.add_landmark_factor}\unskip
, adding the last term of (9.6) as a \pythoninline{BearingRangeFactor2D}. Mind the argument order.
        \begin{longhint}
            The factor is constructed as
            \begin{center}\pythoninline{gtsam.BearingRangeFactor2D(pose_key, landmark_key, bearing, range, noise_model)}\end{center}
            with the bearing \emph{first}, as a \pythoninline{gtsam.Rot2} (\pythoninline{gtsam.Rot2(angle)}), and the range as a float. A plain float bearing raises a \pythoninline{TypeError}. The noise model you are handed is already ordered [bearing, range] to match.
        \end{longhint}
        \begin{solution} \begin{pythoncode}
def add_landmark_factor(
    graph: gtsam.NonlinearFactorGraph,
    pose_key: int,
    landmark_key: int,
    measurement: np.ndarray,
    noise_model: gtsam.noiseModel.Base,
) -> None:
    measured_range, measured_bearing = float(measurement[0]), float(measurement[1])
    graph.add(
        gtsam.BearingRangeFactor2D(
            pose_key,
            landmark_key,
            gtsam.Rot2(measured_bearing),
            measured_range,
            noise_model,
        )
    )
\end{pythoncode}
 \end{solution}
    \end{subparts}

    % ------------------------------------------------------------------ (d)
    \breakans
    \part Implement \pythoninline{factor_graph.predict_measurement}\unskip
, the measurement function \(\mathbf{h}(x_k, m^j)\) of (9.6), together with the two blocks of the Jacobian \(\mathbf{J}\) in (9.8) belonging to this pose and this landmark. It returns the predicted measurement \(z = [\text{range}, \text{bearing}]\), the pose block \(\partial z / \partial x_k\) (\(2\times3\)) and the landmark block \(\partial z / \partial m^j\) (\(2\times2\)). GTSAM computes all of them for you, one row at a time:
    \begin{enumerate}
        \item Preallocate four sub-Jacobians, one row each: range and bearing, each with respect to the pose (\(1\times3\)) and to the landmark (\(1\times2\)), e.g.\ \pythoninline{H_range_pose = np.zeros((1, 3), order="F")}.
        \item Pass them to the two GTSAM calls. Each returns its value and fills in its two sub-Jacobians:
        \begin{center}
            \pythoninline{pose.range(landmark, H_range_pose, H_range_landmark)} \\
            \pythoninline{pose.bearing(landmark, H_bearing_pose, H_bearing_landmark)}
        \end{center}
        \item Stack the rows with \pythoninline{np.vstack}, range first and bearing second, to match the order of \(z\).
    \end{enumerate}

    The prediction does not update the estimate; the factors from (c) do that. It predicts what the local map should look like from \(x_k\), so that the front-end can gate the actual measurements against it in (f) and JCBB.

    \begin{longhint}
        \pythoninline{bearing} returns a \pythoninline{gtsam.Rot2}, not a float; take \pythoninline{.theta()}. The sub-Jacobians must be \pythoninline{float64} with exactly these shapes, or GTSAM raises a \pythoninline{TypeError}. Both are in the tangent space at the pose, the same convention as the pose covariance you recover in (e), so do not replace them with hand-derived \(\partial/\partial[x, y, \psi]\) Jacobians.
    \end{longhint}
    \begin{solution} \begin{pythoncode}
def predict_measurement(
    pose: gtsam.Pose2,
    landmark: np.ndarray,
) -> tuple[np.ndarray, np.ndarray, np.ndarray]:
    landmark = gtsam.Point2(float(landmark[0]), float(landmark[1]))

    H_range_pose = np.zeros((1, 3), order="F")
    H_range_landmark = np.zeros((1, 2), order="F")
    H_bearing_pose = np.zeros((1, 3), order="F")
    H_bearing_landmark = np.zeros((1, 2), order="F")

    predicted_range = pose.range(landmark, H_range_pose, H_range_landmark)
    predicted_bearing = pose.bearing(landmark, H_bearing_pose, H_bearing_landmark).theta()

    z = np.array([predicted_range, predicted_bearing])
    H_pose = np.vstack((H_range_pose, H_bearing_pose))
    H_landmark = np.vstack((H_range_landmark, H_bearing_landmark))

    return z, H_pose, H_landmark
\end{pythoncode}
 \end{solution}

    % ------------------------------------------------------------------ (e)
    \breakans
    \part Implement \pythoninline{factor_graph.assemble_joint_covariance}\unskip
, which builds the joint covariance \(P\) of the current pose and the landmarks in the local map: the \(P\) in \(S = HPH^T + R\) of (f).

    A factor graph stores information, not covariance, so \(P\) has to be recovered from the graph (Sec.\ 9.4). GTSAM does that for you and hands you the result as a \pythoninline{gtsam.JointMarginal}, which stores \(P\) in blocks: \pythoninline{joint_marginal.at(key_i, key_j)} is the covariance between two variables. Your job is to put these blocks together into one matrix, in the order of \pythoninline{keys} -- the pose first, then the landmarks. Either stack the blocks with \pythoninline{np.block}, or preallocate \(P\) and fill each block in two loops.

    Two details matter. Include every block, also those between two landmarks: these cross-covariances are what let JCBB judge a set of associations jointly. And keep the order of \pythoninline{keys}, because (f) stacks the Jacobians in that order. A \(P\) in another order still gives a plausible-looking \(S\), just a wrong one.

    \begin{solution} \begin{pythoncode}
def assemble_joint_covariance(
    joint_marginal: gtsam.JointMarginal,
    keys: list[int],
) -> np.ndarray:
    return np.block([
        [joint_marginal.at(key_i, key_j) for key_j in keys]
        for key_i in keys
    ])

    # Equivalent, with two loops filling a preallocated P:
    #
    # dims = [joint_marginal.at(key, key).shape[0] for key in keys]
    # offsets = np.concatenate([[0], np.cumsum(dims)])  # where each block starts
    # P = np.zeros((offsets[-1], offsets[-1]))
    # for i, key_i in enumerate(keys):
    #     for j, key_j in enumerate(keys):
    #         rows = slice(offsets[i], offsets[i + 1])
    #         cols = slice(offsets[j], offsets[j + 1])
    #         P[rows, cols] = joint_marginal.at(key_i, key_j)
    # return P
\end{pythoncode}
 \end{solution}

    % ------------------------------------------------------------------ (f)
    \breakans
    \part Implement \pythoninline{factor_graph.innovation_covariance}\unskip
, which assembles
    \[
        S = H P H^T + R
    \]
    for the whole local map at once. This is (9.26) with \(P = (R^TR)^{-1}\) recovered in part (e), \(H\) the stacked Jacobian \(\mathbf{J}_a\) from part (d), and \(R\) the block-diagonal measurement noise \(\mathbf{R}_a\).

    \(H\) is dense in its three leftmost columns, because every measurement depends on the pose, and block diagonal elsewhere, because measurement \(i\) depends only on landmark \(i\).

    Note that \(P\) contains the cross-covariances \emph{between landmarks}. Those cross terms are the entire reason JCBB can be joint rather than a sequence of independent gates: they are what makes a set of individually plausible associations collectively implausible.

    \begin{solution} \begin{pythoncode}
def innovation_covariance(
    jacobians_pose: list[np.ndarray],
    jacobians_landmark: list[np.ndarray],
    joint_covariance: np.ndarray,
    measurement_cov: np.ndarray,
) -> np.ndarray:
    n = len(jacobians_pose)

    H = np.zeros((2 * n, 3 + 2 * n))
    for i in range(n):
        H[2 * i : 2 * i + 2, 0:3] = jacobians_pose[i]
        H[2 * i : 2 * i + 2, 3 + 2 * i : 3 + 2 * i + 2] = jacobians_landmark[i]

    R = np.kron(np.eye(n), measurement_cov)

    return H @ joint_covariance @ H.T + R
\end{pythoncode}
 \end{solution}

    % ------------------------------------------------------------------ (g)
    \breakans
    \part Landmark birth. A measurement that JCBB could not associate is either a landmark you have not seen before, or clutter. Putting every one of them straight into the graph fills the map with spurious landmarks that then attract wrong associations; never putting them in means the map never grows.
    \begin{subparts}
        \subpart Implement \pythoninline{factor_graph.inverse_measurement}\unskip
, the inverse measurement function \(\mathbf{h}^{-1}(z, x)\).

        Only the position is needed, not a covariance: once the landmark is connected by factors, the graph works out its uncertainty for itself. That is a direct consequence of (9.2), the posterior is a product of factors. EKF-SLAM (Sec.\ 7.2) has to propagate a new landmark's covariance by hand instead.
        \begin{solution} \begin{pythoncode}
def inverse_measurement(pose: gtsam.Pose2, measurement: np.ndarray) -> np.ndarray:
    measured_range, measured_bearing = float(measurement[0]), float(measurement[1])
    landmark_body = gtsam.Rot2(measured_bearing).rotate(gtsam.Point2(measured_range, 0.0))
    landmark_world = pose.transformFrom(landmark_body)
    return np.asarray(landmark_world, dtype=float).reshape(2)
\end{pythoncode}
 \end{solution}

        \subpart Implement \pythoninline{landmark_manager.TentativeLandmark.is_confirmed}\unskip
, the \(M\) of \(N\) confirmation rule: promote a tentative landmark once it has been seen in at least \(M\) distinct time steps inside a sliding window of the last \(N\), inclusive at both ends. This is the same track initiation logic used in target tracking. Count \emph{time steps}, not observations: here they coincide, because association is one-to-one, but a front-end that allowed two measurements of one landmark in a single scan would otherwise confirm on one scan's worth of evidence -- exactly the clutter burst the rule exists to reject.
        \begin{solution} \begin{pythoncode}
def is_confirmed(self, current_step: int, M: int, N: int) -> bool:
    window_start = current_step - N + 1
    hits_in_window = sum(
        observation.step >= window_start
        for observation in self.supporting_observations
    )
    return hits_in_window >= M
\end{pythoncode}
 \end{solution}
    \end{subparts}

\end{parts}
